% ============================================================
% Juan_Cano_Master_Inventory.tex
% MASTER CAREER INVENTORY — NOT AN APPLICATION DOCUMENT
% No page limit. All bullets, all metrics, all context.
% Source of truth for all tailored derivatives.
% Engine: pdfLaTeX | Font: Helvet
% Last updated: March 2026
% ============================================================
\documentclass[10pt,a4paper]{article}

\usepackage[T1]{fontenc}
\usepackage[utf8]{inputenc}
\usepackage{helvet}
\renewcommand{\familydefault}{\sfdefault}
\usepackage[top=1.8cm, bottom=1.8cm, left=2.0cm, right=2.0cm]{geometry}
\usepackage{titlesec}
\usepackage{enumitem}
\usepackage{hyperref}
\usepackage{microtype}
\usepackage{parskip}
\usepackage{xcolor}
\usepackage{mdframed}

\pagestyle{plain}
\setlength{\parindent}{0pt}
\setlength{\parskip}{2pt}

% --- Section heading style ---
\titleformat{\section}
  {\normalfont\large\bfseries\uppercase}
  {}{0em}{}
  [\vspace{-4pt}\rule{\linewidth}{0.8pt}]
\titlespacing{\section}{0pt}{12pt}{5pt}

\titleformat{\subsection}
  {\normalfont\normalsize\bfseries}
  {}{0em}{}
\titlespacing{\subsection}{0pt}{8pt}{3pt}

% --- List style ---
\setlist[itemize]{
  noitemsep,
  topsep=3pt,
  parsep=0pt,
  partopsep=0pt,
  leftmargin=1.4em,
  label=\textbullet
}

\setlist[itemize,2]{
  noitemsep,
  topsep=1pt,
  leftmargin=1.4em,
  label=--
}

% --- Hyperlink style ---
\hypersetup{colorlinks=true, urlcolor=blue, hidelinks}

% --- Notice box ---
\newmdenv[
  linecolor=black!40,
  backgroundcolor=black!5,
  innerleftmargin=8pt,
  innerrightmargin=8pt,
  innertopmargin=6pt,
  innerbottommargin=6pt,
  skipabove=6pt,
  skipbelow=6pt
]{noticebox}

% --- Role entry ---
\newcommand{\roleentry}[3]{%
  \vspace{4pt}%
  \textbf{#1} \hfill \textit{#2}\\%
  \textit{\small #3}\\[-2pt]%
}

\begin{document}

% ============================================================
% NOTICE
% ============================================================
\begin{noticebox}
\textbf{MASTER INVENTORY — DO NOT SUBMIT DIRECTLY.}
This document is a complete, unfiltered career repository. It contains all roles, all bullet variants,
all metrics, all context notes, and all inventory items including side projects and partial initiatives.
Application-specific baselines (\texttt{Juan\_Cano\_Resume.tex}, \texttt{\_Photo\_EN}, \texttt{\_Photo\_DE})
are distilled from this file. One-page limit applies to those derivatives, not here.
\end{noticebox}

\vspace{6pt}

% ============================================================
% IDENTITY
% ============================================================
{\LARGE\bfseries JUAN CANO-SELLES}\\[2pt]
{\normalsize Technical Program Manager $|$ Senior Product Owner $|$ Cloud \& AI Strategy}\\[4pt]
{\small
  Braunschweig, Lower Saxony, Germany\\
  +49 163 150 9581 \quad\textbar\quad
  \href{mailto:juan.cano.selles@gmail.com}{juan.cano.selles@gmail.com} \quad\textbar\quad
  \href{https://linkedin.com/in/jcanoselles}{linkedin.com/in/jcanoselles}\\[2pt]
  \textbf{German: C1} \quad\textbar\quad
  \textbf{English: C2} \quad\textbar\quad
  Spanish: Native \quad\textbar\quad French: B1\\[2pt]
  Nationality: Spanish — EU citizen, no work permit required\\
  Availability: 2 months notice
}

% ============================================================
% CAREER SUMMARY (raw — rewrite per application)
% ============================================================
\section{Career Summary (raw — rewrite per application)}

10 years of professional experience across systems integration, technical project management, and
senior product ownership in cloud and connectivity. Last 7 years focused on automotive cloud
infrastructure at scale — governing a 5-team Agile Release Train for a 10M+ vehicle platform.
Hands-on background in embedded systems, firmware, and Linux (Akkodis/SEAT). Strong cross-functional
delivery track record in complex, multi-stakeholder environments.

Certifications: PMP (2025), SAFe PO/PM (2021), AWS SAA-C03 (2026), AWS MLA-C01 (in progress).

Target roles: Senior Technical Program Manager, Senior Product Owner, Senior Lead Project Manager.
Target domains: cloud infrastructure, AI/ML programme governance, platform transformation, tech-forward
companies in automotive-adjacent, defense tech, fintech, or enterprise software.

Career principle: will not accept roles with accountability that lack corresponding authority/mandate.

% ============================================================
% PROFESSIONAL EXPERIENCE
% ============================================================
\section{Professional Experience}

% ----------------------------------------------------------
% SABBATICAL
% ----------------------------------------------------------
\roleentry{Sabbatical \& Professional Development}
  {Sep 2025 – Present}
  {Braunschweig, DE}

\textit{Context:} Structured exit from CARIAD SE on mutual agreement (Aufhebungsvertrag, 06.05.2025,
effective 31.08.2025). Not a layoff — voluntary exit under CARIAD restructuring 2023–2025.
Zeugnis grade 2 (gut), clean, ``auf eigenen Wunsch''. Abfindung received Jan 2026.

\begin{itemize}
  \item Completed three months of planned travel (sabbatical phase).
  \item Completed PMP certification (PMI, 2025) — intensive exam preparation via Udemy mock simulator.
  \item Completed AWS Solutions Architect Associate SAA-C03 (2026) — phased study plan, Udemy course + practice tests.
  \item AWS Machine Learning Engineer Associate MLA-C01 in progress (as of March 2026).
  \item Conducted active job search targeting Munich area and Switzerland (CH: Senior PO/TPM salary anchor CHF 140k gross/year).
\end{itemize}

\textit{Aufhebungsvertrag restriction (Ziffer 4):} Cannot rejoin any VW Group entity (§15 AktG)
on a sozialversicherungspflichtig basis before 31.08.2028 without repaying Abfindung pro-rata.
Blocked employers: CARIAD, MOIA, MHP, Porsche Consulting, Audi, SEAT, Škoda, VW Nutzfahrzeuge,
VW Financial Services, all Rivian/VW JVs. Residual liability approx. EUR 112k as of March 2026,
decreasing by 1/36 per month.

\vspace{8pt}

% ----------------------------------------------------------
% CARIAD SE
% ----------------------------------------------------------
\roleentry{CARIAD SE (Volkswagen Group Company) — Senior Product Owner}
  {Nov 2020 – Aug 2025}
  {Wolfsburg, DE}

\textit{Context:} CARIAD SE is the VW Group's software company, created to consolidate automotive
software across all brands (VW, Audi, SEAT, Škoda, Porsche, Lamborghini, Bentley). Juanjo held
the title ``Senior Product Owner'' on paper; the scope and accountability were de facto
Technical Program Manager level — governing a full Agile Release Train, not a single product backlog.
CARIAD went through major public restructuring 2023–2025 (headcount reductions, leadership changes,
program scope reductions). This context is documented and explains the ``schwierigen Arbeitsbedingungen''
referenced in the Zeugnis.

For AI-forward applications, title variant: ``Technical Program Manager $|$ AI Data Strategy \& Cloud Governance''.

\subsection*{Programme Governance \& Platform Scale}

\begin{itemize}
  \item Governed the technical roadmap for an Agile Release Train (ART) comprising five cross-functional
        development teams, ensuring architectural alignment and regulatory compliance across all VW brands
        and markets for a platform serving 10M+ vehicles globally.
  \item Owned backlog prioritisation, PI Planning, and dependency resolution across the ART,
        maintaining strategic coherence between engineering constraints and brand-level business requirements.
  \item Directed quarterly roadmap reviews with senior leadership, translating programme status,
        risks, and delivery commitments into executive-level communication.
  \item Managed cross-brand stakeholder relationships — SEAT, Audi, VW Passenger Cars, Škoda —
        resolving competing priorities across product lines with different market timelines and regulatory landscapes.
  \item Operated in an environment of sustained organisational ambiguity (restructuring 2023–2025),
        maintaining programme delivery continuity while team composition and leadership changed repeatedly.
\end{itemize}

\subsection*{ML \& AI Initiatives (verified, on-paper)}

\begin{itemize}
  \item \textbf{ML-based behavioral anomaly detection (main initiative):} Governed the ART-level
        programme for an ML-based system detecting behavioral anomalies across the connected vehicle fleet.
        Scope included infrastructure delivery, cross-team alignment with data engineering and ML teams,
        and integration into the fleet reliability pipeline. Programme governance role — not ML
        engineering. Fleet scope: 10M+ vehicles.
  \item \textbf{RAG-based LLM pilot (side project, internal transformation team):} Piloted a
        Retrieval-Augmented Generation LLM for departmental onboarding within the internal transformation
        team. Defined requirements, coordinated implementation, evaluated adoption metrics.
        Framing: AI adoption programme management, not ML engineering.
  \item \textbf{PMO automation tools evaluation (side project, internal transformation team):}
        Evaluated PMO automation tooling to reduce high-cost administrative overhead. Assessed
        build-vs-buy, cost/efficiency tradeoffs, and implementation feasibility.
\end{itemize}

\subsection*{Cloud Architecture \& Cost Governance}

\begin{itemize}
  \item Led the architectural transformation from a monolithic legacy system to a microservice-based
        environment, implementing asynchronous communication via Apache Kafka to scale global data
        ingestion and processing across all connected brands.
  \item Enforced a FinOps-aligned cost-optimisation strategy that delivered a 30\% year-on-year
        reduction in AWS cloud expenditure through workload migration and refined data access patterns.
        \textit{(Only verified quantified cost metric — do not fabricate others.)}
  \item Governed the cloud migration programme from Azure to AWS, including workload prioritisation,
        dependency sequencing, and compliance validation across brands.
  \item Directed the infrastructure programme for Functions-on-Demand (FoD) — over-the-air software
        feature activation on 4M+ vehicles — including backend platform, billing integration,
        and B2B fleet monetisation capability.
  \item Managed data residency compliance requirements for China market operations, coordinating
        with legal, security, and engineering teams to meet local regulatory constraints.
\end{itemize}

\subsection*{Connectivity \& Platform Features}

\begin{itemize}
  \item Governed the programme for real-time vehicle telemetry data pipelines (event-driven,
        Kafka-based), enabling fleet analytics and connected services for OEM customers.
  \item Directed backlog for connectivity platform features across multiple vehicle series and
        brand deployments — managing competing feature requests from SEAT, Audi, and VW
        Passenger Cars within a shared platform constraint.
  \item Coordinated integration of MQTT-based communication layers within the vehicle connectivity
        stack, bridging embedded systems domain knowledge with cloud platform delivery.
\end{itemize}

\vspace{8pt}

% ----------------------------------------------------------
% SEAT S.A.
% ----------------------------------------------------------
\roleentry{SEAT S.A. — Technical Project Lead (Connectivity Systems)}
  {Mar 2019 – Nov 2020}
  {Barcelona, ES}

\textit{Context:} SEAT S.A. is the Spanish brand of the VW Group. Role was embedded in the
connectivity engineering team, managing Tier-1 supplier relationships and cross-functional
delivery for production vehicle programs.

\textit{Aufhebungsvertrag note:} SEAT S.A. is a VW Group verbundenes Unternehmen (§15 AktG).
Cannot return before 31.08.2028 without Abfindung repayment.

\begin{itemize}
  \item Managed the technical development and integration of connectivity units for three major
        vehicle series: SEAT Tarraco (SUV), SEAT León (4th gen), and Cupra Born (EV).
  \item Directed cross-functional deliveries between Tier-1 suppliers (Harman, Continental)
        and internal SEAT departments, covering technical documentation, wiring schematics,
        calibration parameters, and hardware/software release cycles.
  \item Enforced compliance with international homologation standards and internal quality gates,
        including EU eCall mandatory certification and safety-critical crash-test preparations.
  \item Coordinated OTA software update validation cycles and managed ECU software versioning
        across production lines.
  \item Managed technical risk in supplier contracts — escalation paths, milestone tracking,
        deviation documentation — under formal project governance.
\end{itemize}

\vspace{8pt}

% ----------------------------------------------------------
% AKKODIS (formerly Modis)
% ----------------------------------------------------------
\roleentry{Akkodis (formerly Modis) — Systems Integration Engineer}
  {Oct 2015 – Mar 2019}
  {Barcelona, ES}

\textit{Context:} On-site resident consultant embedded at SEAT S.A. for the full tenure.
Akkodis is an engineering consulting firm (Adecco Group). Role combined embedded systems
engineering with technical project coordination for series-production connectivity programs.

\begin{itemize}
  \item Steered technical project coordination and supplier integration for series-production
        connectivity units, serving as on-site resident consultant at SEAT S.A.
  \item Architected low-level firmware for NFC-based IoT wearables, optimising power consumption
        and resolving real-time execution constraints for battery-powered devices.
  \item Developed custom Linux drivers and hardware emulators (Python/Qt) to automate
        communication protocol validation and accelerate software testing workflows.
  \item Engineered UART/SPI/I2C protocol validation tooling for embedded connectivity modules,
        reducing manual test overhead in the ECU integration pipeline.
  \item Coordinated cross-functional deliveries between hardware, firmware, and software teams
        during ECU bring-up phases for production connectivity units.
\end{itemize}

% ============================================================
% TECHNICAL SKILLS — FULL INVENTORY
% ============================================================
\section{Technical Skills — Full Inventory}

\subsection*{Programme \& Project Management}
\begin{itemize}
  \item PMP (Project Management Professional) — PMI, 2025
  \item SAFe PO/PM (Scaled Agile Framework) — 2021
  \item Lean Portfolio Management
  \item FinOps — cloud cost governance, workload optimisation, data access pattern refinement
  \item Agile/DevOps Governance — PI Planning, ART facilitation, dependency management
  \item Roadmap development — quarterly, annual, multi-year
  \item Risk management — risk register, mitigation planning, escalation paths
  \item Stakeholder communication — executive reporting, cross-brand alignment
  \item OKR definition and tracking
\end{itemize}

\subsection*{Cloud \& Architecture}
\begin{itemize}
  \item AWS Solutions Architect Associate SAA-C03 — 2026
  \item AWS Machine Learning Engineer Associate MLA-C01 — in progress (March 2026)
  \item Cloud Migration (Azure to AWS) — programme governance, workload sequencing
  \item Event-Driven Architecture — Apache Kafka, asynchronous messaging, data pipelines
  \item Microservice architecture — decomposition strategy, service boundary definition
  \item Data residency compliance — China regulations, cross-border data governance
  \item MQTT — vehicle connectivity protocol, IoT messaging
  \item Kubernetes — conceptual (Coursera, 2025); not hands-on deployment
\end{itemize}

\subsection*{AI \& Machine Learning (programme governance framing)}
\begin{itemize}
  \item ML programme governance — behavioural anomaly detection initiative (CARIAD)
  \item RAG-based LLM implementation — requirements definition, pilot coordination (CARIAD)
  \item PMO automation tooling evaluation — build-vs-buy, cost/efficiency (CARIAD)
  \item AWS MLA-C01 in progress — ML lifecycle, model evaluation concepts, data annotation
        (theoretical foundation, not hands-on ML engineering)
  \item AI for Everyone (Coursera, 2025) — completed
  \item Generative AI fundamentals (Coursera, 2025) — completed
\end{itemize}

\subsection*{Embedded Systems \& Low-Level (Akkodis/SEAT era)}
\begin{itemize}
  \item Firmware development — C, RTOS, NFC, IoT wearables
  \item Linux driver development — custom kernel modules
  \item Hardware emulation — Python/Qt, protocol simulators
  \item Communication protocols — UART, SPI, I2C, CAN, MQTT
  \item ECU integration — bring-up, validation, calibration
  \item Python — scripting, test automation, hardware emulators
\end{itemize}

\subsection*{Tools \& Methodologies}
\begin{itemize}
  \item Jira, Confluence — backlog management, programme documentation
  \item Azure DevOps — sprint tracking, CI/CD pipeline visibility
  \item MS Project — programme scheduling (older roles)
  \item SAFe tooling — PI Planning boards, ART metrics
  \item FinOps tooling — AWS Cost Explorer, billing dashboards
\end{itemize}

\subsection*{Languages}
\begin{itemize}
  \item Spanish — Native
  \item English — C2 (native level), professional and technical
  \item German — C1 (professional proficiency), ongoing active learning
  \item French — B1 (intermediate)
\end{itemize}

% ============================================================
% EDUCATION
% ============================================================
\section{Education}

\roleentry{Master of Science in Telecommunications Engineering}
  {2013 – 2015}
  {Miguel Hernández University (UMH), Elche, Spain}

\begin{itemize}
  \item Specialisation: embedded systems, signal processing, communications networks.
  \item Thesis: low-power NFC system design for wearable IoT devices (direct predecessor to Akkodis firmware work).
\end{itemize}

% ============================================================
% CERTIFICATIONS — FULL LIST
% ============================================================
\section{Certifications}

\begin{itemize}
  \item \textbf{PMP} — Project Management Professional, PMI, 2025
  \item \textbf{AWS SAA-C03} — Solutions Architect Associate, Amazon Web Services, 2026
  \item \textbf{AWS MLA-C01} — Machine Learning Engineer Associate, Amazon Web Services, \textit{in progress} (March 2026)
  \item \textbf{SAFe PO/PM} — Scaled Agile Framework, 2021
\end{itemize}

% ============================================================
% PSYCHOMETRIC / ASSESSMENT RESULTS (for reference)
% ============================================================
\section{Assessment Results (internal reference only)}

\textit{BMW Group — ELIGO cognitive assessment (2025, position closed before outcome):}
\begin{itemize}
  \item Applied Reasoning: 97th percentile (exceptional)
  \item Diagram Analysis: 79th percentile (high)
  \item Problem-Solving Capacity: 71st percentile (medium-high)
  \item Processing Speed: 54th percentile (median — symbol-based time-pressure tasks)
\end{itemize}

% ============================================================
% JOB SEARCH CONTEXT (operational notes)
% ============================================================
\section{Job Search Context — March 2026}

\subsection*{Market \& Location}
\begin{itemize}
  \item Primary market: Munich area (DE) — active, 3 callbacks from 4 applications within 2 weeks.
  \item Secondary market: Switzerland (CH) — active, Zurich/Zug primary targets.
  \item CH salary anchor: CHF 140k gross/year (conservative, 25th–75th percentile CHF 130k–173k for Senior PO/TPM).
  \item CH application requirements: photo, tailored cover letter, Arbeitszeugnis, EU citizen line mandatory.
  \item DE application requirements: photo, Anschreiben, Arbeitszeugnis + diplomas.
\end{itemize}

\subsection*{Target Employers (confirmed watchlist)}
\begin{itemize}
  \item Rheinmetall AG — watchlist (employer target due to German rearmament context; current posting below 75\% fit due to defense/space domain gap).
  \item Apple (Munich) — active application in progress (EPM Annotation \& Evaluation, March 2026).
\end{itemize}

\subsection*{Blocked Employers (Aufhebungsvertrag Ziffer 4)}
\begin{itemize}
  \item All VW Group entities (§15 AktG) until 31.08.2028.
  \item Explicit list: CARIAD, MOIA, MHP, Porsche Consulting, Audi, SEAT, Škoda,
        VW Nutzfahrzeuge, VW Financial Services, all Rivian/VW JV entities.
  \item Residual liability: approx. EUR 112k brutto as of March 2026, decreasing by 1/36 per month.
\end{itemize}

\subsection*{Application Rules}
\begin{itemize}
  \item Minimum fit threshold for active application: 75\%.
  \item Below 75\%: watchlist or skip. Exceptions only with direct internal contact.
  \item Resume: one-page target for CH/DE. Never submit master inventory or baselines directly.
  \item Cover letter: 3 paragraphs, one page strict. No rhetorical flourishes. Acknowledge domain gaps directly.
  \item Scorecard threshold: overall 7.5+, no dimension below 6.
\end{itemize}

\end{document}
